% =====================================================================
%  main.tex  --  "director de orquesta" de la memoria
% =====================================================================
%  Este fichero SOLO ensambla el proyecto. No contiene diseño (está en
%  styles/mistilo.sty) ni texto largo (está en las carpetas):
%     frontmatter/  resumen, abstract y agradecimientos
%     capitulos/    los seis capítulos
%     apendices/    los apéndices
%     bib/refs.bib  la bibliografía
%     Recursos/     imágenes y logo
%
%  CÓMO SE COMPILA (la bibliografía necesita el programa biber):
%     pdflatex main -> biber main -> pdflatex main -> pdflatex main
%  ¿Por qué tantas pasadas? La 1.ª recoge las citas, biber prepara la
%  bibliografía y las dos últimas colocan números, índices y referencias.
%  Con un solo comando:  latexmk -pdf main.tex   (en Overleaf: Recompile)
%
%  CÓMO LEER UN COMANDO:  \comando[opciones]{argumento}
%  Todo lo que sigue a un  %  hasta final de línea es comentario.
% =====================================================================

% Clase "report": estructura de TFG con \chapter. [11pt] = tamaño de letra.
\documentclass[11pt,a4paper]{report}

% ---- Idioma y codificación ----
% babel: el ÚLTIMO idioma de la lista es el principal (spanish). El inglés
% se carga porque el Abstract usa \begin{otherlanguage}{english}.
\usepackage[utf8]{inputenc}
\usepackage[T1]{fontenc}
\usepackage[english,spanish]{babel}

% ---- Bibliografía ----
% backend=biber -> programa que procesa el .bib (necesario).
% style=ieee    -> formato [1], [2]...
% QUÉ PUEDES CAMBIAR: style=authoryear, style=numeric, style=apa...
% Recompila: el texto de los capítulos NO se toca.
\usepackage[backend=biber,style=ieee]{biblatex}
\addbibresource{bib/refs.bib}

% ---- Nuestro paquete de estilo ----
% Ruta relativa a este fichero: el proyecto es portable.
\usepackage{styles/mistilo}

% ---- Datos del documento ----
% Cada TFG nuevo solo cambia estas líneas.
%   \titulodoc[título corto para la cabecera]{título completo}
\titulodoc[Consulta paralela sobre datos enlazados]{Diseño de un prototipo de consulta paralela sobre datos enlazados comprimidos}
\autordoc{Nombre Apellido}                                  % <- tu nombre
\tutordoc{Nombre Apellido}                                  % <- tu tutor/a
\titulacion{Ingeniería informática}                         % <- tu grado
\mencion{Tecnologías de la información y la comunicación}   % <- tu mención
% \fechadoc{Octubre de 2026}   % descomenta para fijar la fecha

% ---- Compilar solo un capítulo mientras trabajas ----
% En una memoria larga, compilar todo tarda. Descomenta para compilar
% únicamente el capítulo 4; los demás conservan su numeración (\include).
% Vuelve a comentarla al terminar o faltarán capítulos en el PDF.
% \includeonly{capitulos/04_implementacion}

\begin{document}

\makeportada                         % portada definida en mistilo.sty

% ---- Páginas previas: numeración romana (i, ii, iii...) ----
\pagenumbering{roman}
% \input inserta el fichero tal cual (sin salto de página extra).
% frontmatter/resumen.tex -- SOLO CONTENIDO
% \capitulosinnumero (definido en mistilo.sty) crea un capítulo sin
% número, lo añade al índice y ajusta la cabecera.
\capitulosinnumero{Resumen}
Este trabajo presenta un prototipo ficticio que reparte consultas sobre
datos enlazados comprimidos entre varios procesos. Se describe su
arquitectura, se implementa en Python y se evalúa con datos ilustrativos.
La memoria sirve como ejemplo de documento largo para practicar la
modularización de proyectos en \LaTeX.

% frontmatter/abstract.tex
% otherlanguage cambia el idioma SOLO dentro del entorno (separación
% silábica y textos automáticos). Requiere cargar english en babel.
\capitulosinnumero{Abstract}
\begin{otherlanguage}{english}
This work presents a fictitious prototype that distributes queries over
compressed linked data among several processes. It is only an example
document used to practise how to split a \LaTeX{} project into modules.
\end{otherlanguage}

% frontmatter/agradecimientos.tex
\capitulosinnumero{Agradecimientos}
A quienes han hecho posible este taller.


\renewcommand{\contentsname}{Índice}
\tableofcontents                     % índice general automático
\listoffigures                       % índice de figuras (usa \caption)
\listoftables                        % índice de tablas

% ---- Cuerpo: numeración normal (1, 2, 3...) ----
\cleardoublepage
\pagenumbering{arabic}

% \include empieza en página nueva y permite \includeonly.
% PARA AÑADIR UN CAPÍTULO NUEVO:
%   1) crea capitulos/07_xxx.tex que empiece por \chapter{...}
%   2) añade aquí:  \include{capitulos/07_xxx}
% =====================================================================
%  capitulos/01_introduccion.tex  --  SOLO CONTENIDO
% =====================================================================
%  Este fichero contiene únicamente texto. Lo que NO hay (a propósito):
%    - \documentclass, \usepackage ni \begin{document}: eso es de main.tex
%    - colores, márgenes ni formatos: eso es diseño y vive en el .sty
%  PISTAS PARA EDITAR
%    - Una línea en blanco separa párrafos.
%    - ~ es un espacio que NO permite salto de línea ("capítulo~2").
%    - \emph{...} = cursiva; \textbf{...} = negrita.
% =====================================================================
% \chapter crea un capítulo numerado y lo añade al índice.
\chapter{Introducción}

% =====================================================================
%  capitulos/02_estado_arte.tex
%  Aquí practicamos: macros del .sty (\HDT, \RDF), notas al pie,
%  figuras, tablas y citas de distintos tipos de entrada.
% =====================================================================
\chapter{Estado del arte}
\label{cap:arte}

% =====================================================================
%  capitulos/03_diseno.tex
%  Aquí practicamos: tablas, figuras, definiciones, ecuaciones y
%  ejemplos. Los entornos definicion y ejemplo están en mistilo.sty.
% =====================================================================
\chapter{Diseño}
\label{cap:diseno}


% =====================================================================
%  capitulos/04_implementacion.tex
%  Aquí practicamos los listados de código. El estilo "python" y el
%  estilo "terminal" están definidos UNA vez en mistilo.sty; en cada
%  listado solo se escribe [style=...].
% =====================================================================
\chapter{Implementación}
\label{cap:impl}


% =====================================================================
%  capitulos/05_resultados.tex
%  Tabla de resultados, gráfica y referencia a la definición del cap. 3.
%  NOTA: los datos son FICTICIOS (así se indica en el texto).
% =====================================================================
\chapter{Resultados}
\label{cap:result}


% capitulos/06_conclusiones.tex
\chapter{Conclusiones}
\label{cap:conc}



% ---- Apéndices ----
% \appendix hace que los capítulos siguientes se numeren A, B, C...
\appendix
% apendices/A_instalacion.tex
% Tras \appendix (en main.tex) este \chapter se numera como "A".
\chapter{Instalación}
\begin{lstlisting}[style=terminal]
$ python -m venv venv
$ source venv/bin/activate
$ pip install -r requirements.txt
\end{lstlisting}

% apendices/B_datos.tex
\chapter{Conjuntos de datos}
% \, es un espacio fino: sirve para separar miles (10\,000).
\begin{table}[H]
  \centering
  \begin{tabular}{@{}lrr@{}}
    \toprule
    \textbf{Nombre} & \textbf{Tripletas} & \textbf{Tamaño (MB)} \\
    \midrule
    pequeno & 10\,000     &  1 \\
    mediano & 1\,000\,000 & 80 \\
    \bottomrule
  \end{tabular}
  \caption{Conjuntos de datos de prueba (ficticios).}
  \label{tab:datos}
\end{table}


% ---- Bibliografía ----
% Imprime solo las referencias realmente citadas con \cite.
\printbibliography[heading=bibintoc,title={Bibliografía}]

\end{document}
